\clearpage
\section[Boundary expectations and reconstruction of the analytic channels]{Boundary expectations and\\reconstruction of the analytic channels}
\label{sec:lectores-gamma}
% Fuente íntegra: ampliacion/expediente/AGENTE_CARTAS_Y_EXPECTATIVA.md

% PROCEDENCIA Fecha del corte: 11 de septiembre de 2026. Nota interna de desarrollo, no modificación de un artículo ni dictamen global sobre HMT. No se ha comunicado a otras tareas del usuario.

\subsection{Action of the readouts and content of the moments}
\label{gamma-add-sub-5}

The explicit actions of the odometer, the difference readout, the uniform inclusion, the Helmert refinement, the bidirectional boundary, and the polar chart have been recovered and composed. The calculation yields:

\begin{enumerate}
\item The moments of the positive and negative columns, considered separately, are determined by explicit operators. They include all cross products between test functions.
\item The uniform expectation does not coincide with the Helmert refinement. On the carry mode it yields an exact factor, calculated below; it does not by itself convert a translation difference into the uniform mode.
\item The gamma-tail resolvent does perform the analytic conversion and exactly reconstructs all negative rows, on every window and before any sign is used.
\item The simultaneous identification of the first moment of the encoder and the moment of its projection does not follow from the three preceding chart identities by a simple direct sum. The specific additional composition that must be verified is identified below. The structural readout is not declared absent, nor is this focal check identified with the full TPK.

\end{enumerate}
The advance consists in calculating the action and cross products, not in reiterating that the projection is positive.

\paragraph{Provenance of the charts.}
The formulas are assembled from the developments \emph{Polar spectral proof},
\emph{Form and Gram}, \emph{Gamma-tail connector}, \emph{Odometer and
carry}, \emph{Pulsation and expectation}, and \emph{Modular polar chart}.
The locators and the record of the reading performed are retained in the
source dossier accompanying the delivery. Each chart is used with the spaces,
normalizations, and identities made explicit below; a
documentary reference does not replace these conditions.

% DOCUMENTACION ## 2. Propietarios y lectura efectuada
% DOCUMENTACION 
% DOCUMENTACION Se utilizan estos propietarios, leídos completos en esta formalización:
% DOCUMENTACION 
% DOCUMENTACION - **P10**: [10b — prueba espectral polar](</Users/ruben/Documents/New project/output/ARTICULO_VIII_PRIMOS_ESTRUCTURA_ESPECTRAL_20260910/antecedentes/integral/manuscrito/sucesor_102/deltas_ley9/10b_prueba_espectral_polar_riemann_weil_rectificacion_probatoria_20260903.tex>), 1.166 líneas. Las columnas están en 151–235; la expectativa y el codificador nativo en 288–367; la transferencia derivada de los momentos en 369–452; el codificador escrito mediante el complemento, más adelante, utiliza la identidad energética ya obtenida.
% DOCUMENTACION - **PG**: [Forma y Gram](</Users/ruben/Documents/excelencia academica/ARTICULOS_CIENTIFICOS_HMT_MD_2026-08-14/RIEMANN_REDHEFFER_WEIL/manuscrito/secciones/00_forma_y_gram.tex>), 630 líneas. Fórmulas explícitas de los dos modos y del refinamiento: 133–191; columnas: 194–220; codificador y momentos comunes: 222–304.
% DOCUMENTACION - **PT**: [Conector de cola gamma](</Users/ruben/Documents/excelencia academica/ARTICULOS_CIENTIFICOS_HMT_MD_2026-08-14/RIEMANN_REDHEFFER_WEIL/manuscrito/secciones/02b_conector_cola_gamma.tex>), completo. Inversión: 1–45; elevación y frontera: 47–155; todas las filas negativas: 157–242; restricción alcanzable: 250–269.
% DOCUMENTACION - **PO**: [Odómetro y acarreo](</Users/ruben/Documents/excelencia academica/CLAUSURA_WEIL_HMT_PULSACION_NONADICA_2026-08-11/PUBLICACION_WEIL_RIEMANN_HMT_MD_REVISION_EDITORIAL_2026-08-12/manuscrito/sections/03_odometro_carry.tex>), completo.
% DOCUMENTACION - **PE**: [Pulsación y expectativa](</Users/ruben/Documents/excelencia academica/CLAUSURA_WEIL_HMT_PULSACION_NONADICA_2026-08-11/PUBLICACION_WEIL_RIEMANN_HMT_MD_REVISION_EDITORIAL_2026-08-12/manuscrito/sections/02_pulsacion_y_expectativa.tex>), completo.
% DOCUMENTACION - **PP**: [Carta polar modular](</Users/ruben/Documents/excelencia academica/CLAUSURA_WEIL_HMT_PULSACION_NONADICA_2026-08-11/PUBLICACION_WEIL_RIEMANN_HMT_MD_REVISION_EDITORIAL_2026-08-12/manuscrito/sections/06_canal_polar_modular.tex>), completo. Carta Witt–dodecafase: 134–181.
% DOCUMENTACION - [Acarreo ortogonal del residual](</Users/ruben/Documents/excelencia academica/ARTICULOS_CIENTIFICOS_HMT_MD_2026-08-14/RIEMANN_REDHEFFER_WEIL/manuscrito/secciones/00d_carry_ortogonal_app_weil.tex>), completo.
% DOCUMENTACION 
% DOCUMENTACION Se contrastó además el propietario de las cinco realizaciones del continuo del 17 de agosto, sección de Riemann desde su objeto hasta el transporte del complemento, sin disgregar esas construcciones: allí se vuelve a formular el mismo par de momentos nativos. La consulta de esa sección no equivale a una lectura íntegra de todo ese manuscrito.
% DOCUMENTACION 
% DOCUMENTACION Los antecedentes de agosto se consultan como propietarios de las fórmulas que P10 reutiliza, no como sustitutos de la selección autoral del integral de 2.249 páginas. Los localizadores globales que todavía resuelven 2.084 páginas quedan registrados como históricos para esta serie.
% DOCUMENTACION 
\subsection{Conventions and genealogical cutoff}
\label{gamma-add-sub-32}

The relevant starting point is the APP--TRIT--TPK state, its refinement with memory, and its native readouts. In this formalization:

\begin{itemize}
\item APP provides the support and the irreducible selector. The native product, coproductivity, and selector precede the representation of prime--power labels.
\item TRIT and sheet transport preserve orientation and the seam.
\item The TPK preserves the phase step and the carry register; its odometer induces the translation difference. The logarithmic representation of the clocks produces \(\ell_{p,k}=k\log p\) and \(w_{p,k}=(\log p)p^{-k/2}\).
\item The boundary and polar charts are subsequent realizations of the same readouts. No state is selected by means of zeta zeros or a target value.

\end{itemize}
This section calculates the arrow between these readouts and the expectation, after their common generation; it does not seek to reconstruct or replace the joint continuum.

We work with orthonormal history coordinates. In these coordinates, \(u_N=N^{-1/2}(1,\ldots,1)\). If functions with the uniform probability measure are used, the constant coordinates and averages are conjugated by the corresponding normalization factor; the projections and all the norm ratios below remain the same.

Let \(H=L^2(\mathbb R)\), \(T_af(x)=f(x-a)\), \(\mathcal D_R=C_c^\infty((-R/2,R/2))\). We write \(d_a=T_a-I\); changing simultaneously to \(I-T_a\) preserves their Grams but changes the sign of the inverting cotrace.

The inner product is linear in the first variable. We retain the
polar evaluations and the scalar of the full form,
\begin{equation}
\begin{aligned}
A_f&=\int_{\mathbb R}f(x)e^{x/2}\,dx,
&B_f&=\int_{\mathbb R}f(x)e^{-x/2}\,dx,\\
c_\Gamma&=-\gamma-\frac\pi2-3\log2-\log\pi<0,
&\mu(a)&=\frac{e^{-a/2}}{1-e^{-2a}}.
\end{aligned}
\label{gamma-add-def-evaluaciones}
\end{equation}
Here \(\pi\), \(\gamma\), and the logarithms retain the internal
representation and subsequent realization declared in the preceding chapters.
The formula introduces no target values into the generator.

\subsection{Action of the odometer on generators}
\label{gamma-add-sub-47}

For \(N=9^m\), \(m\ge1\), with the seam located at the first vector,

\begin{equation}
U_{a,N}(e_r\otimes f)=
\begin{cases}
e_{r+1}\otimes f,&r<N-1,\\
e_0\otimes T_af,&r=N-1.
\end{cases}
\label{gamma-add-display-1}
\end{equation}

The uniformization \(i_Nf=u_N\otimes f\) gives

\begin{equation}
U_{a,N}i_Nf
=u_N\otimes f+N^{-1/2}e_0\otimes d_af.
\label{gamma-add-display-2}
\end{equation}

Therefore, for \(P_N=|u_N\rangle\langle u_N|\otimes I\),

\begin{equation}
\begin{aligned}
i_N^*U_{a,N}i_N&=I+\frac{d_a}{N},\\
(I-P_N)U_{a,N}i_N
&=\frac{\sqrt{N-1}}{N}\eta_N\otimes d_a.
\end{aligned}
\label{gamma-add-display-3}
\end{equation}

where

\begin{equation}
\eta_N=\frac{Ne_0-\mathbf1_N}{\sqrt{N(N-1)}}.
\label{gamma-add-display-4}
\end{equation}

This recovers the normalized difference readout:

\begin{equation}
\widehat\delta_{a,m}f=\eta_N\otimes u_{81}\otimes d_af,
\qquad
j_mf=u_N\otimes u_{81}\otimes f.
\label{gamma-add-display-5}
\end{equation}

The action on any combination of generators, without removing their cross terms, is

\begin{equation}
\left\langle\sum_i\widehat\delta_{a_i,m}f_i,
                  \sum_j\widehat\delta_{b_j,m}g_j\right\rangle
=\sum_{i,j}\langle d_{a_i}f_i,d_{b_j}g_j\rangle.
\label{gamma-add-display-6}
\end{equation}

In particular,

\begin{equation}
\widehat\delta_{a,m}^*\widehat\delta_{b,m}=d_a^*d_b,\qquad
j_m^*j_m=I,\qquad
j_m^*\widehat\delta_{a,m}=0.
\label{gamma-add-eq-1}
\end{equation}

The last cross product follows from \(\langle u_N,\eta_N\rangle=0\). It is concrete information about the orientation of the modes, in addition to the two isolated norms.

\subsection{Two distinct refinements and their expectations}
\label{gamma-add-sub-108}

\subsubsection{Uniform inclusion and projection of the last digit}
\label{gamma-add-sub-110}

Write \(M=9N\) and \(\mathbb C^M=\mathbb C^N\otimes\mathbb C^9\). The inclusion and its adjoint are

\begin{equation}
\iota_Nx=x\otimes u_9,\qquad
E_N=I_N\otimes\langle u_9|,\qquad P_N^{\rm unif}=\iota_NE_N.
\label{gamma-add-display-8}
\end{equation}

Acting in coordinates,

\begin{equation}
E_Nu_M=u_N,\qquad
E_N\eta_M=a_N\eta_N,\qquad
a_N=\sqrt{\frac{N-1}{9N-1}}.
\label{gamma-add-eq-2}
\end{equation}

\begin{proof}
We have \(E_Ne_0^{(M)}=e_0^{(N)}/3\) and \(E_N\mathbf1_M=3\mathbf1_N\).
Substitution into \(\eta_M=(Me_0-\mathbf1_M)/\sqrt{M(M-1)}\)
yields \eqref{gamma-add-eq-2}.
\end{proof}

Consequently,

\begin{equation}
(E_N\otimes I)\widehat\delta_{a,m+1}=a_N\widehat\delta_{a,m},
\qquad
(E_N\otimes I)j_{m+1}=j_m.
\label{gamma-add-eq-3}
\end{equation}

All projected moments are determined:

\begin{equation}
\begin{aligned}
\widehat\delta_{a,m+1}^*(P_N^{\rm unif}\otimes I)\widehat\delta_{b,m+1}
 &=a_N^2d_a^*d_b,\\
j_{m+1}^*(P_N^{\rm unif}\otimes I)j_{m+1}&=I,\\
j_{m+1}^*(P_N^{\rm unif}\otimes I)\widehat\delta_{b,m+1}&=0.
\end{aligned}
\label{gamma-add-eq-4}
\end{equation}

The complementary fraction of the carry is \(1-a_N^2=8N/(9N-1)\). For \(N=9\), the expressed fraction is \(1/10\), not \(1\), and its output remains a difference.

We do not conflate \(P_N^{\rm unif}\), which averages the last digit, with the projection onto the single global line \(u_M\): the latter annihilates \(\eta_M\) entirely.

\subsubsection{Helmert refinement}
\label{gamma-add-sub-156}

The operator \(W_N\) is fixed by the orthonormal Helmert bases
of the two levels, ordered to preserve the uniform and
carry modes. It is the isometry that sends the first \(N\) vectors of the
basis at level \(N\) to the corresponding vectors of the basis at level \(M\).
In particular, it satisfies

\begin{equation}
W_Nu_N=u_M,\qquad W_N\eta_N=\eta_M,\qquad W_N^*W_N=I.
\label{gamma-add-display-12}
\end{equation}

Thus, \(E_N^{\rm H}=W_N^*\), \(P_N^{\rm H}=W_NW_N^*\) give

\begin{equation}
E_N^{\rm H}j_{m+1}=j_m,\qquad
E_N^{\rm H}\widehat\delta_{a,m+1}=\widehat\delta_{a,m}.
\label{gamma-add-eq-5}
\end{equation}

The uniform--carry cross product remains zero, and that of two differences remains \(d_a^*d_b\). The refinement is therefore natural for both modes, but does not transform one into the other. Specifically, \(W_N\ne\iota_N\).

\subsubsection{Conjugation of expectations and preservation of moments}
\label{gamma-add-sub-174}

There is a unitary \(S_M\) on \(\mathbb C^M\) that maps \(\iota_N\mathbb C^N\) to \(W_N\mathbb C^N\), with \(W_N=S_M\iota_N\). It suffices to map the two orthonormal bases to each other and complete the complements. Then

\begin{equation}
P_N^{\rm H}=S_MP_N^{\rm unif}S_M^*.
\label{gamma-add-display-14}
\end{equation}

Conjugating the state \(C\) and its projection simultaneously,

\begin{equation}
(S_MC)^*(S_MP S_M^*)(S_MC)=C^*PC,\qquad
(S_MC)^*(S_MC)=C^*C.
\label{gamma-add-eq-6}
\end{equation}

Neither moment changes. If only the projection is changed, the fraction in this case goes from \(a_N^2\) to \(1\), but the analytic operator remains \(d_a\), not \(I\).

More generally, for any projection acting only on the finite history fiber, the moment of a single channel is
\(\langle\eta,P\eta\rangle d_a^*d_b\). This change of basis does not turn it into \(2I\). Any effective coupling between analytic channels is a different operation and must be written down, not inferred from conjugation.

\subsection{Metric identities of the boundary and polar charts}
\label{gamma-add-sub-195}

\subsubsection{Bidirectional boundary}
\label{gamma-add-sub-197}

The bidirectional chart preserves the composition

\begin{equation}
\begin{aligned}
K_2&=\mathbb C^{81}\otimes u_9^\perp,\\
\Delta_a&=j_{\rm APP}(\eta_9\otimes d_a),\\
\beta_0&=R_{\rm tw}(L^3+\varepsilon\nabla_x)FI_{Q_9}|_{K_2}.
\end{aligned}
\label{gamma-add-display-16}
\end{equation}

In this realization, \(j_{\rm APP}\) is the isometric inclusion
of the indicated fiber; \(\beta_0\) is injective, and
\(\Lambda_{\rm tw}^{\#}\) is positive definite on the boundary
space. These are the conditions of the chart: the factors of the
composition are supplied with compatible domains, and that
injectivity is not deduced here from an equality of dimensions.
% ENSAMBLAJE: los factores R_tw, L^3, epsilon*nabla_x, F e I_Q9
% pertenecen al propietario PT, 85--155. Para esta identidad bastan
% la inclusión isométrica, beta0 inyectivo y Lambda_tw^# positivo.
With \(V_\beta=\beta_0(\beta_0^*\beta_0)^{-1/2}\) and
\(U_{\rm tw}=(\Lambda_{\rm tw}^{\#})^{-1/2}V_\beta\),

\begin{equation}
G_a=(U_{\rm tw}\otimes I)\Delta_a,\qquad
G_a^*(\Lambda_{\rm tw}^{\#}\otimes I)G_b=d_a^*d_b.
\label{gamma-add-eq-7}
\end{equation}

\begin{proof}
The polar normalization gives \(V_\beta^*V_\beta=I\) and therefore
\(U_{\rm tw}^*\Lambda_{\rm tw}^{\#}U_{\rm tw}=I\).
Since \(j_{\rm APP}\) and \(\eta_9\) are normalized,
\(\Delta_a^*\Delta_b=d_a^*d_b\). Composition proves
\eqref{gamma-add-eq-7}, including all cross products.
\end{proof}
\(U_{\rm tw}\) is isometric for the declared metric
\(\Lambda_{\rm tw}^{\#}\). In ordinary Hilbert-space coordinates,
the isometry between the effective spaces is
\((\Lambda_{\rm tw}^{\#})^{1/2}U_{\rm tw}=V_\beta\).

Identity \eqref{gamma-add-eq-7} contains neither an expectation nor its value on \(G_a\). To transport a projection, one must conjugate it with that same isometry and its adjoint in the correct metric. When this is done, \eqref{gamma-add-eq-6} preserves the preceding compressed moment: passing to the boundary does not create the second moment.

\subsubsection{Effective polar chart}
\label{gamma-add-sub-220}

The polar realization by five octads supplies the five-column
matrix \(T\), whose Gram and normalization are

\begin{equation}
T^*T=4I_5+\frac43J_5,\qquad
\widehat T=T(T^*T)^{-1/2},\qquad
F_{34}=(f_3,f_6,f_9,f_4,f_8).
\label{gamma-add-display-18}
\end{equation}

The chart, with its spaces made explicit, is

\begin{equation}
U=F_{34}\widehat T^*:
\operatorname{Ran}\widehat T\subset\mathbb C^{24}
\longrightarrow\operatorname{span}(f_3,f_6,f_9,f_4,f_8)\subset\mathbb C^{12}.
\label{gamma-add-display-19}
\end{equation}

For \(\eta=\operatorname{diag}(1,1,1,-1,-1)\),

\begin{equation}
U^*J_{34}U=\widehat T\eta\widehat T^*,\qquad
\widehat T^*U^*J_{34}U\widehat T=\eta.
\label{gamma-add-eq-8}
\end{equation}

Here the five vectors of \(F_{34}\) are orthonormal, and
\(F_{34}^*J_{34}F_{34}=\eta\). Substituting
\(U=F_{34}\widehat T^*\) and \(\widehat T^*\widehat T=I_5\)
proves both equalities in \eqref{gamma-add-eq-8}.
This is the typed formulation of the equality with the diagonal matrix.
The five roles are effective; the identification of coordinates must
not be erased.
% ENSAMBLAJE: la selección concreta de cinco octadas y el cálculo de
% T*T=4I5+(4/3)J5 proceden de PP, 134--181. Se conserva esta carta con
% su Gram de partida explícito; no se genera aquí un nuevo diseño.

For the two polar evaluations,

\begin{equation}
b_\pm(f)=\frac{A_f\pm B_f}{\sqrt2},\qquad
b_+(f)\overline{b_+(g)}-b_-(f)\overline{b_-(g)}
=A_f\overline{B_g}+B_f\overline{A_g}.
\label{gamma-add-eq-9}
\end{equation}

\eqref{gamma-add-eq-8} transports the signature, and \eqref{gamma-add-eq-9} transports its cross terms. Neither identifies the isolated positive line as a sufficient source of the negative one: for \(f\) real, odd, and nonzero, chosen with \(f(x)>0\) on the positive half-axis of its support, we have \(b_+(f)=0\) and \(b_-(f)>0\). Recovery must involve the common analytic state, as does the gamma-tail connector developed below.

\subsection{Difference between channelwise readout and common transfer}
\label{gamma-add-sub-259}

To fix all the summands used, let
\(\mathcal K_m=\mathbb C^{9^m}\otimes\mathbb C^{81}\otimes H\),
and let \(\zeta_+\), \(\zeta_-\) be unit vectors of the two
polar lines. The indices \(\ell\leq R\) range over the prime--power
labels \((p,k)\) with \(k\log p\leq R\), preserving
their weight \(w_\ell=(\log p)p^{-k/2}\). The columns are
\begin{equation}
\begin{aligned}
J_{+,R}f={}&\Bigl(
 (\sqrt{w_\ell}\,\widehat\delta_{\ell,m}f)_{\ell\leq R},\,
 (\sqrt{\mu(a)}\,\widehat\delta_{a,m}f)_{a>0},\,
 \zeta_+b_+(f)\Bigr),\\
J_{-,R}f={}&\Bigl(
 \sqrt{-c_\Gamma}\,j_mf,\,
 (\sqrt{2w_\ell}\,j_mf)_{\ell\leq R},\,
 \zeta_-b_-(f)\Bigr).
\end{aligned}
\label{gamma-add-columnas-completas}
\end{equation}
The continuous component of the first column belongs to
\(L^2((0,\infty),da;\mathcal K_m)\); the remaining components
are assembled in the indicated orthogonal sums. The evaluations
\(b_\pm\) are those of \eqref{gamma-add-eq-9}. Their common domain
consists of test functions with finite gamma energy. The identity
\begin{equation}
\mathscr W_R(f,g)
=\langle J_{+,R}f,J_{+,R}g\rangle
 -\langle J_{-,R}f,J_{-,R}g\rangle
\label{gamma-add-diferencia-columnas}
\end{equation}
is verified by expanding
\(\langle d_\ell f,d_\ell g\rangle
=2\langle f,g\rangle-\langle T_\ell f,g\rangle
-\langle f,T_\ell g\rangle\), using the gamma integral
and the polar identity \eqref{gamma-add-eq-9}. This preserves
the three channels and all cross terms.

Taking the explicit positive column \(\Xi_+\) as the encoder satisfies its first moment by direct calculation. Subsequently applying the uniform expectation produces, in each prime channel,

\begin{equation}
w_\ell a_N^2\langle d_\ell f,d_\ell g\rangle,
\label{gamma-add-display-22}
\end{equation}

whereas the required negative row is

\begin{equation}
2w_\ell\langle f,g\rangle.
\label{gamma-add-display-23}
\end{equation}

These are not the same identity. For example, choose \(R>\ell>0\) and \(f\in\mathcal D_R\) real and nonnegative, with support overlapping its translation by \(\ell\). Then

\begin{equation}
\|d_\ell f\|^2
=2\|f\|^2-2\langle f,T_\ell f\rangle
<2\|f\|^2.
\label{gamma-add-eq-10}
\end{equation}

No contractive projection applied solely to that channel can produce \(\sqrt2f\). The Helmert refinement removes the factor \(a_N^2\), but \eqref{gamma-add-eq-10} persists. This is a falsifier of the \textbf{channel-by-channel} identification, not a refutation of a common transfer that redistributes energy among the channels.

On the other hand, directly introducing \(j_mf\), its prime--scalar copies, and \(b_-(f)\) as orthogonal summands of the encoder ensures their outputs but adds their norms to the first moment. Preserving the original first moment requires an explicit isometric redistribution. That redistribution is precisely the operation that the simultaneous identification of the two moments must exhibit.

Invariance of the cyclic subspace under \(P\) does not by itself determine the value of \(PCf\) either; it only guarantees that this value remains in the subspace.

\subsection{Gamma-tail inverse and reconstruction of all rows}
\label{gamma-add-sub-288}

Here there is indeed an explicit action on every generator and every linear combination. Let

\begin{equation}
\mu(a)=\frac{e^{-a/2}}{1-e^{-2a}},\quad
A_R=[R+1,R+2],\quad \nu_R=\int_{A_R}\mu(a)\,da>0,
\quad M_R=\mathbf1_{(-R/2,R/2)}.
\label{gamma-add-display-25}
\end{equation}

For the sign \(d_a=T_a-I\) already fixed, define the cotrace

\begin{equation}
\epsilon_m=\langle\eta_{9^m}|\otimes\langle u_{81}|\otimes M_R,
\qquad
\mathcal R^-_{m,R}F
=-\nu_R^{-1}\int_{A_R}\sqrt{\mu(a)}\,\epsilon_m F(a)\,da.
\label{gamma-add-eq-11}
\end{equation}

The minus sign disappears if \(I-T_a\) is used, which is the alternative convention for the tail readout. For \(f\in\mathcal D_R\) and \(a>R\), \(M_RT_af=0\). Consequently,

\begin{equation}
\mathcal R^-_{m,R}
\bigl(a\mapsto\sqrt{\mu(a)}\widehat\delta_{a,m}f\bigr)=f.
\label{gamma-add-eq-12}
\end{equation}
\label{gamma-add:cola-lector}

Moreover, \(\|\mathcal R^-_{m,R}\|\le\nu_R^{-1/2}\) by Cauchy--Schwarz. No individual translation is inverted, and Weil positivity is not used.

Write

\begin{equation}
S_Rf=\left(\sqrt{-c_\Gamma}\,j_mf,\,
              (\sqrt{2w_\ell}\,j_mf)_{\ell\le R}\right),
\quad B_-f=\zeta_-b_-(f).
\label{gamma-add-display-28}
\end{equation}

The concrete common row, defined on the full positive ambient space, is

\begin{equation}
\mathcal C_R^{\rm tail}
=
\begin{pmatrix}
0&S_R\mathcal R^-_{m,R}&0\\
0&B_-\mathcal R^-_{m,R}&0
\end{pmatrix}.
\label{gamma-add-eq-13}
\end{equation}

From \eqref{gamma-add-eq-12},

\begin{equation}
\mathcal C_R^{\rm tail}J_{+,R}f=J_{-,R}f
\quad\hbox{for every }f\in\mathcal D_R.
\label{gamma-add-eq-14}
\end{equation}

This action contains the scalar, prime--power, and polar rows, with their constants and orientations intact. For any pair \(f,g\),

\begin{equation}
\begin{aligned}
\langle\mathcal C_R^{\rm tail}J_+f,\mathcal C_R^{\rm tail}J_+g\rangle
={}&\left(-c_\Gamma+2\sum_{\ell\le R}w_\ell\right)\langle f,g\rangle\\
&+b_-(f)\overline{b_-(g)}.
\end{aligned}
\label{gamma-add-eq-15}
\end{equation}

For \(f=\sum_i f_i\), \(g=\sum_j g_j\), the right-hand side contains the full sum over \(i,j\); the test functions have not been replaced by orthogonal cells. The prime--gamma cross terms of \eqref{gamma-add-eq-7} are preserved when the input representations are changed. The final form is

\begin{equation}
J_+^*\left(I-(\mathcal C_R^{\rm tail})^*\mathcal C_R^{\rm tail}\right)J_+
=J_+^*J_+-J_-^*J_-=\mathcal W_R.
\label{gamma-add-eq-16}
\end{equation}

This is an effective equality of forms, without yet assuming that the operator in parentheses is positive. The operator is bounded; an explicit bound is

\begin{equation}
\|\mathcal C_R^{\rm tail}\|
\le\nu_R^{-1/2}
\sqrt{-c_\Gamma+2\sum_{\ell\le R}w_\ell+2\sinh(R/2)-R}.
\label{gamma-add-display-33}
\end{equation}

The bound fixes the type; it does not prove that the norm is at most one.

\subsection{Exact condition for combining both moments in one expectation}
\label{gamma-add-sub-375}

Let \(\mathcal M_+=\overline{J_+\mathcal D_R}\). The reachable transfer

\begin{equation}
C_{\rm rea}(J_+f)=J_-f
\label{gamma-add-display-34}
\end{equation}

is well defined because the gamma tail proves that \(J_+\) is injective. It is the restriction of \eqref{gamma-add-eq-13} and is therefore bounded. Any other bounded transfer giving the same outputs coincides with it on \(\mathcal M_+\).

A change of coordinates does not provide a different reachable transfer. The condition to be established for \eqref{gamma-add-eq-13} to have an isometric dilation with the native expectation is

\begin{equation}
\|C_{\rm rea}x\|\le\|x\|,
\qquad x\in\mathcal M_+.
\label{gamma-add-eq-17}
\end{equation}

Its relation to the two moments can be stated unambiguously. \textbf{If} \eqref{gamma-add-eq-17} is proved from the native operators, \(D=(I-C_{\rm rea}^*C_{\rm rea})^{1/2}\) exists, and

\begin{equation}
\mathfrak C_{\rm dil}f
=u_9\otimes(J_-f,0)+\eta_9\otimes(0,DJ_+f)
\label{gamma-add-eq-18}
\end{equation}

satisfies, for \(P=|u_9\rangle\langle u_9|\otimes I\),

\begin{equation}
\mathfrak C_{\rm dil}^*\mathfrak C_{\rm dil}=J_+^*J_+,\qquad
\mathfrak C_{\rm dil}^*P\mathfrak C_{\rm dil}=J_-^*J_-.
\label{gamma-add-display-37}
\end{equation}

\eqref{gamma-add-eq-18} is a normal form of the dilation once \eqref{gamma-add-eq-17} has been proved, \textbf{not} a new prospective proof of \eqref{gamma-add-eq-17}. It identifies exactly what the native colligation must produce: the defect of that common transfer, with all its cross terms, and not merely the defect of each clock separately.

The specific next composition combines the energy budget of the full connector with the recovered local and refinement identities. If an alternative native colligation realizes \eqref{gamma-add-eq-14} and is contractive on all of \(\mathcal M_+\), reachable uniqueness proves \eqref{gamma-add-eq-17}. There is no need to invent another readout; its action and common energy do need to be calculated.

\subsection{Checks and provenance}
\label{gamma-add-sub-412}

The actions \eqref{gamma-add-eq-1}, \eqref{gamma-add-eq-7},
\eqref{gamma-add-eq-8}, and \eqref{gamma-add-eq-12}--\eqref{gamma-add-eq-16}
assemble pre-existing operators and results. The compositions
\eqref{gamma-add-eq-2}--\eqref{gamma-add-eq-6} and the falsifier
\eqref{gamma-add-eq-10} are explicit calculations on these formulas;
no mathematical priority is claimed for them, nor is any
global absence of other operators inferred from them.

We retained 25 rational checks on
\(v_{9N}=9Ne_0-\mathbf1\), for \(N=2,3,9,81,729\).
They check the action of the average, the fraction \((N-1)/(9N-1)\),
the zero mean, the orthogonality of the detail, and the complement
\(8N/(9N-1)\). The 25 checks gave the same result in normal
and optimized modes, without assertions that can be disabled. They verify the identities
\eqref{gamma-add-eq-2}--\eqref{gamma-add-eq-4}, not global positivity.
The programs, commands, receipts, and locators are retained in
Appendix~\ref{ap:procedencia-reproduccion}.

% PROCEDENCIA HISTORICA \begin{itemize}
% PROCEDENCIA HISTORICA \item Las acciones \eqref{gamma-add-eq-1}, \eqref{gamma-add-eq-7}, \eqref{gamma-add-eq-8}, \eqref{gamma-add-eq-12}--\eqref{gamma-add-eq-16} reúnen resultados y operadores preexistentes. Estatuto: \textbf{RESULTADO_RECUPERADO / FORMALIZACIÓN_REUNIDA}.
% PROCEDENCIA HISTORICA \item Las composiciones explícitas \eqref{gamma-add-eq-2}--\eqref{gamma-add-eq-6} y el control \eqref{gamma-add-eq-10} son cálculos de esta formalización sobre esas fórmulas; no se les atribuye prioridad matemática ni se los presenta como ausencia global de otros operadores. Procedencia exhaustiva de la forma exacta de esos cálculos: no investigada.
% PROCEDENCIA HISTORICA \item Se ejecutaron 25 controles racionales sobre las coordenadas sin normalizar \(v_{9N}=9Ne_0-\mathbf1\), para \(N=2,3,9,81,729\). Verificaron la acción del promedio, la fracción \((N-1)/(9N-1)\), la media nula, la ortogonalidad del detalle y la fracción complementaria \(8N/(9N-1)\). Se conservaron en el script reproducible, sin uso de aserciones eliminables. Pasaron los mismos 25 controles en modo normal y los 25 en modo optimizado: recibo normal y recibo optimizado. Son controles de \eqref{gamma-add-eq-2}--\eqref{gamma-add-eq-4}, no de positividad global.
% PROCEDENCIA HISTORICA \item Pasaron el núcleo formal permanente, la autocomprobación causal de constantes, el comprobador del resolutor y la puerta de principios operativos. El resolutor conserva la referencia histórica de 2.084 páginas; esta formalización respeta la selección autoral de 2.249 y no modifica ese registro.
% PROCEDENCIA HISTORICA \item El expediente fuente conserva los PDF, índices y paquetes anteriores inalterados; esta incorporación añade la exposición de las identidades y no cambia retroactivamente sus recibos.
% PROCEDENCIA HISTORICA 
% PROCEDENCIA HISTORICA \end{itemize}
% PROCEDENCIA HISTORICA Para repetir los controles desde esta carpeta:
% PROCEDENCIA HISTORICA 
% PROCEDENCIA HISTORICA % DOCUMENTACION     python3 -I -S verificar_expectativa_modos.py --receipt CONTROL_EXPECTATIVA_MODOS_NORMAL.json
% PROCEDENCIA HISTORICA % DOCUMENTACION     python3 -I -S -O verificar_expectativa_modos.py --receipt CONTROL_EXPECTATIVA_MODOS_OPTIMIZADO.json
% PROCEDENCIA HISTORICA 
% PROCEDENCIA HISTORICA 
\textbf{Operational conclusion:} the state and the two readouts have not disappeared. The calculation recovers their action and demonstrates precisely the difference between isometric transport of modes, uniform expectation, and analytic recovery of channels. The second moment identity has a concrete verification route through the common connector; conjugation of charts alone does not carry out that verification.
